\documentclass[UTF8,zihao=-4,openany]{ctexbook}

\usepackage[a4paper,top=25mm,bottom=25mm,left=25mm,right=25mm]{geometry}
\usepackage{xcolor}
\usepackage{array}
\usepackage{booktabs}
\usepackage{pdfpages}
\usepackage{hyperref}
\usepackage{bookmark}

\definecolor{accent}{HTML}{2F5D73}
\hypersetup{
  unicode=true,
  colorlinks=false,
  pdfborder={0 0 0},
  bookmarksopen=true,
  bookmarksnumbered=false,
  pdftitle={高中物理竞赛一轮讲义：2020春电磁学全十讲},
  pdfauthor={高中物理竞赛一轮讲义}
}
\setlength{\parindent}{0pt}
\pagestyle{empty}

\begin{document}

\pdfbookmark[0]{标题页}{cover}
\begin{titlepage}
  \centering
  \vspace*{40mm}
  {\color{accent}\rule{0.72\textwidth}{1.1pt}}\par
  \vspace{13mm}
  {\zihao{1}\bfseries 高中物理竞赛一轮讲义\par}
  \vspace{9mm}
  {\zihao{2}\bfseries 2020春电磁学全十讲\par}
  \vspace{8mm}
  {\zihao{-2}第32—41节\quad 含解答与解图\par}
  \vspace{13mm}
  {\color{accent}\rule{0.72\textwidth}{1.1pt}}\par
  \vfill
  {\zihao{-4}依据《2020春高一物理》第1—10讲整理\par}
  \vspace{14mm}
\end{titlepage}

\clearpage
\pdfbookmark[0]{目录}{contents}
\begin{center}
  {\zihao{2}\bfseries 目录}\par
  \vspace{10mm}
\end{center}

\renewcommand{\arraystretch}{1.42}
\begin{tabular}{>{\raggedright\arraybackslash}p{0.13\textwidth}
                >{\raggedright\arraybackslash}p{0.67\textwidth}
                >{\raggedleft\arraybackslash}p{0.08\textwidth}}
  \toprule
  \textbf{讲次} & \textbf{主题} & \textbf{PDF页} \\
  \midrule
  第1讲（第32节） & 电场的基本性质 & 3 \\
  第2讲（第33节） & 电势能与电势 & 15 \\
  第3讲（第34节） & 静电平衡 & 30 \\
  第4讲（第35节） & 电介质、电容器 & 46 \\
  第5讲（第36节） & 电流 & 63 \\
  第6讲（第37节） & 复杂电阻网络 & 79 \\
  第7讲（第38节） & Lorentz 力 & 94 \\
  第8讲（第39节） & Ampere 力，Ampere 定律 & 115 \\
  第9讲（第40节） & 电磁感应：基本原理 & 132 \\
  第10讲（第41节） & 补充阅读：电磁场与超导体（选讲） & 157 \\
  \bottomrule
\end{tabular}

\vfill
{\small\color{gray}目录页码为合订本 PDF 的物理页码。}

\includepdf[
  pages=-,
  pagecommand={},
  addtotoc={1,chapter,1,{第1讲：电场的基本性质},{lecture:1}}
]{../thirty-second-section/build/main.pdf}

\includepdf[
  pages=-,
  pagecommand={},
  addtotoc={1,chapter,1,{第2讲：电势能与电势},{lecture:2}}
]{../thirty-third-section/build/main.pdf}

\includepdf[
  pages=-,
  pagecommand={},
  addtotoc={1,chapter,1,{第3讲：静电平衡},{lecture:3}}
]{../thirty-fourth-section/build/main.pdf}

\includepdf[
  pages=-,
  pagecommand={},
  addtotoc={1,chapter,1,{第4讲：电介质、电容器},{lecture:4}}
]{../thirty-fifth-section/build/main.pdf}

\includepdf[
  pages=-,
  pagecommand={},
  addtotoc={1,chapter,1,{第5讲：电流},{lecture:5}}
]{../thirty-sixth-section/build/main.pdf}

\includepdf[
  pages=-,
  pagecommand={},
  addtotoc={1,chapter,1,{第6讲：复杂电阻网络},{lecture:6}}
]{../thirty-seventh-section/main.pdf}

\includepdf[
  pages=-,
  pagecommand={},
  addtotoc={1,chapter,1,{第7讲：Lorentz 力},{lecture:7}}
]{../thirty-eighth-section/build/main.pdf}

\includepdf[
  pages=-,
  pagecommand={},
  addtotoc={1,chapter,1,{第8讲：Ampere 力，Ampere 定律},{lecture:8}}
]{../thirty-ninth-section/build/main.pdf}

\includepdf[
  pages=-,
  pagecommand={},
  addtotoc={1,chapter,1,{第9讲：电磁感应：基本原理},{lecture:9}}
]{../fortieth-section/build/main.pdf}

\includepdf[
  pages=-,
  pagecommand={},
  addtotoc={1,chapter,1,{第10讲：补充阅读——电磁场与超导体（选讲）},{lecture:10}}
]{../forty-first-section/build/main.pdf}

\end{document}
